\documentclass[conference]{IEEEtran}

\usepackage[T1]{fontenc}
\usepackage[utf8]{inputenc}
\usepackage{cite}
\usepackage{amsmath,amssymb,amsfonts}
\usepackage{graphicx}
\usepackage{booktabs}
\usepackage{multirow}
\usepackage{array}
\usepackage{url}
\usepackage{xcolor}
\usepackage{enumitem}
\usepackage{longtable}
\usepackage{threeparttable}

\title{NEXUS-ATMS: AI-Centric Urban Congestion Intelligence from Simulation to Pseudo-Live Operations\\
\large IEEEtran Paper Setup Template (Structure + Audit Blueprint)}

\author{%
\IEEEauthorblockN{Author Name 1, Author Name 2, Author Name 3}
\IEEEauthorblockA{Department, Institution, City, Country \\
Email: author@domain.com}
}

\begin{document}
\maketitle

\begin{abstract}
% This is a setup placeholder, not the final abstract text.
% Fill in 150--250 words with: problem context, gap, AI contribution,
% methodology variants, key metrics, and headline results.
% Example anchors from project artifacts:
% - RL vs baseline waiting-time reduction
% - queue-length reduction
% - LSTM forecasting quality
% - anomaly detector F1 and recall
\textbf{Template note:} This file defines what the paper should contain end-to-end,
centered on AI components and full system auditability, without writing the final narrative.
\end{abstract}

\begin{IEEEkeywords}
Intelligent Transportation Systems, Reinforcement Learning, Traffic Signal Control,
Computer Vision, SUMO, FastAPI, Digital Twin, Time-Series Forecasting, Anomaly Detection,
Explainable AI
\end{IEEEkeywords}

\section{Introduction}
\subsection{Urban Congestion Problem Context}
% Include: congestion impact, delay, emissions, emergency response, scalability challenges.

\subsection{Research Gap}
% Explain why isolated optimization (only simulation, only CV, or only static control)
% is insufficient for adaptive city-scale operations.

\subsection{AI-Centric Contributions}
% Convert to final numbered contributions in paper draft.
\begin{enumerate}[leftmargin=*,label=\textbf{C\arabic*.}]
	\item Unified AI stack integrating RL control, CV sensing, prediction, anomaly detection, and specialty modules.
	\item Dual-mode pipeline: SUMO simulation and pseudo-live dashboard mode.
	\item Auditable control loop with REST/WebSocket telemetry, replay, and decision traces.
	\item Comparative evaluation against fixed-timing baseline and cross-module performance analysis.
\end{enumerate}

\subsection{Paper Roadmap}
% Briefly map sections in final version.

\section{Related Work}
\subsection{RL for Traffic Signal Control}
% DQN/PPO/A2C literature and multi-objective signal optimization.

\subsection{Simulation-Driven ITS using SUMO}
% Prior SUMO + TraCI ecosystems and transfer-to-real limitations.

\subsection{Computer Vision for Traffic Sensing}
% YOLO-based counting/tracking/speed estimation and fallback pipelines.

\subsection{Forecasting and Anomaly Detection in Traffic Systems}
% LSTM/Seq2Seq predictors, statistical and ML anomaly detectors.

\subsection{Gaps Relative to This Work}
% Position your novelty: integrated modular operations + live control surface + auditability.

\section{System Scope and Audit Baseline}
\subsection{Project Snapshot}
% Add commit hash/date, hardware profile, software versions.

\subsection{Execution Modes}
\begin{itemize}[leftmargin=*]
	\item \textbf{Simulation/Training Mode:} SUMO-driven RL training and evaluation.
	\item \textbf{Dashboard Demo/Pseudo-Live Mode:} FastAPI + WebSocket streaming with CV and fallback handling.
\end{itemize}

\subsection{End-to-End Data and Control Flow}
% Describe the pipeline from sensing to action to dashboard visualization.
\begin{enumerate}[leftmargin=*]
	\item Ingestion: camera/SUMO/IoT signals.
	\item AI processing: detection, fusion, prediction, anomaly analysis.
	\item Control: RL/manual/emergency phase decisions.
	\item Serving: FastAPI endpoints and WebSocket system state.
	\item Visualization: map, KPIs, incidents, replay timeline.
\end{enumerate}

\section{Dataset Description and Data Sources}
\subsection{Simulation Assets (SUMO)}
% Mention network and route assets; include junction topology and scenario generation.
\begin{itemize}[leftmargin=*]
	\item Single intersection network and routes.
	\item Grid networks (2x2 and configured 4x4 assets).
	\item Scenario route definitions: rush-hour, normal, night, asymmetric.
\end{itemize}

\subsection{Pseudo-Live and Sensor Data Sources}
\begin{itemize}[leftmargin=*]
	\item Camera stream source (OpenCV capture).
	\item Sensor simulator streams (loop, radar, environmental, pedestrian, emergency).
	\item Optional MQTT publication and in-process bus fallback.
\end{itemize}

\subsection{Derived Data Products}
\begin{itemize}[leftmargin=*]
	\item Training logs and checkpoints.
	\item Evaluation JSON and visual reports.
	\item Anomaly and LSTM result artifacts.
	\item Replay/audit logs from dashboard backend.
\end{itemize}

\subsection{Data Splits and Temporal Protocol}
% Specify train/val/test windows, episode split, and leakage controls.

\section{Feature Engineering}
\subsection{RL State Representation}
% Include state dimensions and signal-phase encoding.
\begin{itemize}[leftmargin=*]
	\item SUMO env state: queue lengths, waiting times, phase one-hot, phase elapsed.
	\item Standalone env state: per-approach queue/wait/arrival/occupancy + contextual flags.
\end{itemize}

\subsection{Vision-Derived Features}
\begin{itemize}[leftmargin=*]
	\item Vehicle detections and class-confidence filtering.
	\item Marker geo-projection and lane-level counts.
	\item Density and congestion level estimates.
\end{itemize}

\subsection{IoT Fusion Features}
\begin{itemize}[leftmargin=*]
	\item Vehicle count, occupancy, speed, queue estimate, flow, pedestrian requests.
	\item Data quality scores and environmental attributes.
\end{itemize}

\subsection{Prediction and Anomaly Feature Sets}
\begin{itemize}[leftmargin=*]
	\item LSTM inputs: approach-wise traffic + temporal cyclic embeddings + weather/emergency flags.
	\item Anomaly detector streams: rolling statistical signals and rate-of-change channels.
\end{itemize}

\section{Preprocessing Pipeline}
\subsection{Simulation Preprocessing}
% Scenario loading, warm-up strategy, step interval normalization.

\subsection{Vision Preprocessing}
% Frame resize, confidence thresholding, backend fallback logic.

\subsection{Sensor and Fusion Preprocessing}
% Fault filtering, weighted smoothing, plausibility checks, missing data handling.

\subsection{Time-Series Preprocessing for LSTM}
% Sliding windows, horizon preparation, normalization strategy, sequence assembly.

\subsection{Quality Guards and Contract Validation}
% Include system-state schema checks and consistency warnings.

\section{Approaches}
\subsection{Baseline Approach}
% Fixed-timing signal control baseline definition.

\subsection{RL Approaches}
% DQN and PPO setup, action semantics, reward weighting, training callbacks.

\subsection{Predictive Approach}
% Seq2Seq/BiLSTM forecasting pipeline and forecasting horizons.

\subsection{Anomaly Detection Approach}
% Statistical and ML ensemble methods with alert quorum design.

\subsection{Operations Approach}
% FastAPI + WebSocket control-plane approach with replay and audit endpoints.

\section{Proposed Methodology I: SUMO-Centric AI Signal Optimization}
\subsection{Problem Formulation}
% State, action, reward, transition assumptions.

\subsection{Training Loop Definition}
% Episode protocol, exploration policy, checkpoint strategy.

\subsection{Evaluation Protocol}
% Metrics, baselines, and reproducibility settings.

\section{Proposed Methodology II: Vision+IoT Pseudo-Live AI Orchestration}
\subsection{Live Runtime Architecture}
% CV detection/tracking, map telemetry, state synthesis, and dashboard broadcast.

\subsection{Control Modes}
\begin{itemize}[leftmargin=*]
	\item AI mode (adaptive phase selection).
	\item Manual override mode.
	\item Emergency corridor mode.
\end{itemize}

\subsection{Safety, Security, and Explainability Hooks}
% Pedestrian safety, cybersecurity checks, decision feeds, audit logging.

\section{Experimental Results and Audit Summary}
\subsection{Canonical Source of Results}
% Use results/evaluation_results.json as the canonical source for the RL comparison table.
% The generated HTML report currently shows a stale throughput entry; do not use it as the
% primary citation source for throughput until it is regenerated.
% Training duration should be reported from the run logs: DQN quick-train with 50,000 timesteps
% completed in approximately 18 minutes 13 seconds for the recorded run.

\subsection{Table of Core RL Results}
\begin{table}[t]
\centering
\caption{Canonical evaluation results from \texttt{results/evaluation\_results.json}.}
\label{tab:rl-results}
\begin{threeparttable}
\begin{tabular}{lccc}
\toprule
\textbf{Metric} & \textbf{Fixed-Timing Baseline} & \textbf{RL Agent} & \textbf{Change} \\
\midrule
Mean reward & $-1149.14 \pm 317.28$ & $-20.66 \pm 2.20$ & $+98.20\%$ \\
Avg waiting time (s) & 581.31 & 10.23 & $-98.24\%$ \\
Avg queue length (veh) & 26.16 & 2.62 & $-89.97\%$ \\
Throughput (veh/hr) & 510.8 & 453.8 & $-11.16\%$ \\
Episode length & 720 & 720 & -- \\
\bottomrule
\end{tabular}
\begin{tablenotes}
\footnotesize
\item \textit{Audit note:} The JSON artifact is treated as canonical because it stores both raw baseline/agent metrics and the computed comparison deltas.
\item \textit{Interpretation:} The RL policy sharply reduces delay and queueing, while sacrificing some throughput to prioritise congestion relief.
\end{tablenotes}
\end{threeparttable}
\end{table}

\subsection{Training Time Summary}
\begin{table}[t]
\centering
\caption{Training and model-building time reported for the current project run.}
\label{tab:training-time}
\begin{tabular}{lccc}
\toprule
\textbf{Model / Run} & \textbf{Mode} & \textbf{Timesteps / Epochs} & \textbf{Training Time} \\
\midrule
DQN quick-train RL policy & SUMO demo mode & 50,000 timesteps & 18 min 13 s \\
LSTM traffic predictor & CPU-based training & 50 epochs & \textasciitilde 2 min \\
\bottomrule
\end{tabular}
\end{table}

\subsection{Table of Forecasting and Anomaly Results}
\begin{table}[t]
\centering
\caption{Prediction and anomaly-detection results already produced by the project.}
\label{tab:ml-results}
\begin{tabular}{lcc}
\toprule
\textbf{Component} & \textbf{Metric(s)} & \textbf{Value(s)} \\
\midrule
LSTM traffic predictor & $R^2$ / MAE & 0.6126 / 0.0746 \\
LSTM traffic predictor & Seq2Seq bi-LSTM setup & 24-step input, 6-step horizon \\
IsolationForest & Accuracy & 0.850 \\
Autoencoder & Accuracy & 0.900 \\
Ensemble anomaly detector & Precision / Recall / F1 & 0.840 / 1.000 / 0.913 \\
\bottomrule
\end{tabular}
\end{table}

\subsection{Runtime and System Health Audit}
\begin{table}[t]
\centering
\caption{Observed runtime behavior from the live dashboard backend.}
\label{tab:runtime-audit}
\begin{tabular}{l p{0.61\linewidth}}
\toprule
\textbf{Item} & \textbf{Observed state} \\
\midrule
Backend mode & FastAPI + WebSocket backend serving the dashboard on port 8000 \\
Vision stack & YOLOv8n loaded on CPU with OpenCV capture active \\
Video input & Camera source 0 connected; first valid frame detected in the live runtime \\
Data stream & `system_state` broadcast over WebSocket with frame ID and tick timestamp \\
Health status & Live runtime reported healthy or degraded depending on frame quality \\
Fallback policy & Synthetic / fallback handling available when camera or detector quality is low \\
Audit trail & Replay buffer and audit log endpoints are present for post-hoc inspection \\
\bottomrule
\end{tabular}
\end{table}

\subsection{Audit Coverage Matrix}
\begin{table*}[t]
\centering
\caption{Audit coverage of the current project, from SUMO to OpenCV and the dashboard control plane.}
\label{tab:audit-coverage}
\begin{tabular}{p{0.14\linewidth} p{0.18\linewidth} p{0.20\linewidth} p{0.42\linewidth}}
\toprule
\textbf{Layer} & \textbf{Component} & \textbf{Evidence in Repository} & \textbf{What the paper should report} \\
\midrule
Simulation & SUMO / TraCI environment & \texttt{src/envs/sumo\_env.py}, network XMLs, scenario YAMLs, \texttt{train.py}, \texttt{evaluate.py} & Single-intersection and grid-based signal control, state/action design, reward shaping, evaluation baseline \\
Standalone AI env & First-principles traffic model & \texttt{control/traffic\_env.py} & Poisson demand, queue evolution, reward terms, emergency/pedestrian hooks \\
Vision & OpenCV + YOLOv8 detector & \texttt{vision/detector.py}, live backend logs & Vehicle detection, tracking-ready outputs, CPU/GPU backend selection, fallback chain \\
IoT fusion & Sensor simulator and fusion & \texttt{iot/sensor\_simulator.py}, \texttt{iot/data\_fusion.py} & Loop/radar/environment/pedestrian/emergency streams, weighted fusion, smoothing \\
Prediction & LSTM forecaster & \texttt{prediction/lstm\_predictor.py}, benchmark results & Short-horizon traffic forecasting, sequence design, R\textsuperscript{2}/MAE \\
Anomaly detection & Statistical + ML ensemble & \texttt{prediction/anomaly\_detector.py}, ML anomaly outputs & Alert quorum, precision/recall/F1, safety relevance \\
RL control & DQN/PPO agents & \texttt{src/agents/}, benchmark results & Baseline-vs-agent comparison, action semantics, reward trade-offs \\
Specialty modules & Emergency, carbon, safety, security, maintenance, NLP, voice, counterfactual & \texttt{modules/} subtree & Multi-module architecture and operator-facing control decisions \\
Backend & FastAPI, REST, WebSocket, replay, audit & \texttt{dashboard/backend/main.py}, status API, replay buffer & Real-time system-state schema, telemetry cadence, auditability \\
Frontend & Map dashboard and controls & \texttt{dashboard/frontend/index.html} & Operator interface, live sync indicators, replay timeline, command buttons \\
Evaluation & Metrics and visual reports & \texttt{results/evaluation\_results.json}, \texttt{docs/benchmarks.md}, \texttt{results/report.html} & Canonical numbers, comparison table, hardware/runtime context \\
\bottomrule
\end{tabular}
\end{table*}

\subsection{How to Frame the Results in the Paper}
\begin{itemize}[leftmargin=*]
	\item Use the RL table as the main quantitative result for traffic control quality, especially the waiting-time reduction.
	\item Report training time separately so the professor can see how long the model took to learn before evaluation.
	\item Use the forecasting/anomaly table to show that the system is not only a controller, but an AI sensing-and-prediction stack.
	\item Use the runtime audit table to show that the platform was actually executed, streamed, and validated.
	\item Use the coverage matrix to prove that the paper spans the full pipeline: SUMO, OpenCV, IoT fusion, RL, backend, and dashboard.
\end{itemize}

\section{Discussion}
\subsection{Interpretation of AI Gains and Trade-offs}
% Example: wait reduction vs throughput trade-off.

\subsection{Generalization and Deployment Considerations}
% Domain shift from simulation to real roads, sensor reliability, and calibration.

\subsection{Failure Modes and Mitigations}
% Include camera dropouts, degraded CV, fallback behavior, and API-security checks.

\section{Conclusion and Future Work}
% Summarize completed capabilities and realistic next steps.

\section*{Appendix A: Full Component Audit Checklist}
\renewcommand{\arraystretch}{1.15}
\begin{longtable}{p{0.18\linewidth} p{0.25\linewidth} p{0.24\linewidth} p{0.27\linewidth}}
\toprule
\textbf{Layer} & \textbf{Component} & \textbf{Artifacts to Cite} & \textbf{What to Report in Paper} \\
\midrule
\endhead
Data/Simulation & SUMO environments & train/evaluate scripts, network and route XMLs, scenario YAMLs & topology, state/action spaces, timing parameters, reward design \\
Data/Simulation & Standalone env & control traffic environment implementation & queue model, demand curves, multi-objective reward terms \\
Vision & Detector/tracker/counting/speed/incidents & vision modules and runtime status fields & backend selection, confidence thresholds, fallback chain, detection quality \\
IoT Fusion & Sensor simulator and fusion & iot simulator/fusion modules & sensor types, noise/fault model, weighting and smoothing strategy \\
RL & DQN/PPO agents & src agents, config hyperparameters, training logs & architecture, policy setup, training budget, callbacks \\
Prediction & LSTM forecasting & prediction LSTM module and result artifacts & feature schema, horizon, seq length, accuracy metrics \\
Anomaly & Statistical + ML anomaly & anomaly modules and anomaly result files & detector ensemble logic, quorum, precision/recall/F1 \\
Specialty Modules & Emergency/carbon/pedestrian/security/maintenance/NL/counterfactual/voice & modules folder by submodule & input-output contracts, role in control loop, constraints \\
Backend/API & FastAPI + WebSocket + audit/replay & dashboard backend and endpoints & system-state schema, control endpoints, timing cadence, auditability \\
Frontend/Twin & Dashboard map, controls, replay, health panels & dashboard frontend assets & operator workflow, visualization semantics, live update strategy \\
Evaluation & Baseline and report generation & evaluate script, results JSON, report artifacts & metric definitions, statistical summaries, comparison table design \\
\bottomrule
\end{longtable}

\section*{Appendix B: Reproducibility and Environment Setup Checklist}
\begin{itemize}[leftmargin=*]
	\item Python version, dependency versions, CUDA availability, GPU model.
	\item SUMO installation and TraCI integration status.
	\item Configuration snapshot (default + scenario overrides).
	\item Random seed policy and deterministic settings.
	\item Exact commands for training, evaluation, and dashboard launch.
	\item Runtime mode flags (live, demo, data-only, fallback policies).
	\item Known limitations and unresolved integration gaps.
\end{itemize}

\section*{Appendix C: Suggested Figure and Table Captions (Draft Stubs)}
\begin{itemize}[leftmargin=*]
	\item \textbf{Fig. 1}: Layered architecture of NEXUS-ATMS from data ingestion to AI decision serving.
	\item \textbf{Fig. 2}: Signal control loop showing sensing, policy decision, enforcement, and dashboard feedback.
	\item \textbf{Fig. 3}: Comparative RL training dynamics and evaluation outcomes.
	\item \textbf{Table I}: Hyperparameters and system runtime configuration used for all experiments.
	\item \textbf{Table II}: Baseline versus AI model performance across congestion KPIs.
\end{itemize}

\section*{How to Use This Setup File}
\begin{enumerate}[leftmargin=*]
	\item Copy this \texttt{.tex} file into Overleaf as your main manuscript starter.
	\item Replace placeholders section-by-section with final narrative and citations.
	\item Attach figures, metric tables, and endpoint/module audit evidence as you finalize experiments.
	\item Keep the appendix audit matrix during drafting to ensure no subsystem is omitted.
\end{enumerate}

\bibliographystyle{IEEEtran}
% Create references.bib in Overleaf and switch this on when ready.
% \bibliography{references}

\end{document}
